\documentclass[tikz,border=10pt]{standalone}

\usepackage[a4paper,left=5.5cm,right=1.5cm,top=5cm,bottom=1.5cm]{geometry}

\usepackage{amsmath,amsthm,amssymb,amsfonts}
% \usepackage{array}
\usepackage{booktabs}
\usepackage[table]{xcolor}
\usepackage{cuted}
\usepackage{xcolor}
\usepackage{graphicx}
\graphicspath{{./img/}}
\usepackage{setspace}
\newcommand*\diff{\mathop{}\!\mathrm{d}}
\newcommand*\Diff[1]{\mathop{}\!\mathrm{d^#1}}
\usepackage{enumitem}

\usepackage[utf8]{inputenc}
\usepackage[T1]{fontenc}
\usepackage[spanish]{babel}


\usepackage{tabularx}
\newcolumntype{Y}{>{\centering\arraybackslash}X}




\usepackage[locale=FR, per-mode=fraction, separate-uncertainty=true]{siunitx}
\usepackage{physics}
\AtBeginDocument{\RenewCommandCopy\qty\SI}

\usepackage{pgfplots}
\usepackage{adjustbox, tikz}
\usepackage{tikz-3dplot}
\usepackage{circuitikz}
% \usepackage{gnuplottex}

% \usepackage{ifpdf}
\ifpdf%
  \usepackage{pdftricks}
  \begin{psinputs}
    \usepackage{pst-optic}
%\usepackage{pstricks-add} 	
  \end{psinputs}
\else
  \usepackage{pst-optic}
  \newenvironment{pdfpic}{}{}
\fi
 

\usetikzlibrary{decorations}
\usetikzlibrary{decorations.pathmorphing}
\usetikzlibrary{shapes.geometric}
\usetikzlibrary{arrows}
\usetikzlibrary{arrows.meta}   %edg
\usetikzlibrary{calc}
\pgfplotsset{compat=newest}
\usetikzlibrary{patterns}

\def\cosa{-0.69466}
\def\sina{-0.71934}
% A custom arrowhead for use on x, y, z axes
\tikzstyle{axisarrow} = [-{Latex[inset=0pt,length=5pt]}]

\begin{document}


% \begin{center}
%     \begin{circuitikz}[scale=1]
%       \draw (0,0) to[battery2, invert, l=$6\,\text{V}$] (0,2)
%       to[battery2, invert, l=$4\,\text{V}$] (0,4)
%       to[R, l_=$30\,\Omega$] (3.5,4) to (5.5,4)
%       to[battery2, l=$6\,\text{V}$] (5.5,0)
%       to (3.5,0)
%       to[R, l_=$50\,\Omega$] (0,0);
%       \draw (0,2) to[R, l=\mbox{$R_0 = 50\,\Omega$}] (3.5,2);
%       \draw (3.5,4) to[R, l=$R_1$] (3.5,2)
%       to[R, l=$R_2$] (3.5,0);

%     %   (0,2.5) to[R=$R$, color=red] (2.5,2.5) to[R=$R$, color=red] (5,2.5)
%     %   (5,0) -- (2.5,0) to[R=$R$, color=red] (0,0)
%     %   (0,0) -- (1,1)
%     %   (1.5,1.5) -- (2.5,2.5);
%     %   \fill (1,1) circle (3pt);
%     %   \fill (1.5,1.5) circle (3pt);
%     %   \draw (1,1) node[right]{$A$};
%     %   \draw (1.5,1.5) node[right]{$B$};
%     \end{circuitikz}
% \end{center}


Calcular la energía almacenada en el capacitor cuando el circuito se encuentra operando en régimen estacionario.

\begin{center}
  \begin{circuitikz}[scale=1]
      \draw (2.5,0) to[battery2, invert, l=$110\,\text{V}$] (2.5,1.5) to[R, l=$14.2\,\text{k}\Omega$] (2.5,3.5);
      \draw (0,0) to[R, l=$10.5\,\text{k}\Omega$] (0,3.5) -- (5,3.5)
      to[R, l_=$22.1\,\text{k}\Omega$] (5,1.5)
      to[C, l_=$8.55\,\text{mF}$] (5,0) -- (0,0);
  \end{circuitikz}
\end{center}


\end{document}




